
\documentclass[12pt,a4paper,onecolumn,oneside]{memoir}
%twoside for print version
\RequirePackage[xetex]{graphicx}
\graphicspath{ {}{img/} }

\usepackage[xetex,table,xcdraw]{xcolor}

\usepackage{tabularx}

\usepackage[
breaklinks=true,colorlinks=true,
linkcolor=blue,urlcolor=blue,citecolor=blue,% PDF VIEW
%linkcolor=black,urlcolor=black,citecolor=black,% PRINT
bookmarks=true,bookmarksopenlevel=2]{hyperref}

\usepackage{geometry}
% PDF VIEW
% \geometry{total={210mm,297mm},
% left=25mm,right=25mm,%
% bindingoffset=0mm, top=25mm,bottom=25mm}
% PRINT
\geometry{total={210mm,297mm},
left=20mm,right=20mm,
bindingoffset=10mm, top=25mm,bottom=25mm}

\OnehalfSpacing
%\linespread{1.3}

\usepackage{cite}

% TABLES
\usepackage{tablefootnote}
\usepackage{booktabs}
\usepackage{pbox}
\renewcommand{\arraystretch}{1.2} 
%\renewcommand*\cmidrule{} % No table middle lines
%\renewcommand{\arraystretch}{1.5} % Additional spacing with no middle lines
%\renewcommand*\cmidrule{\hdashline[1pt/2pt]}% Dashed middle lines
\renewcommand*\cmidrule{\midrule[0.001em]} % Thin table middle lines
%\renewcommand*\cmidrule{\midrule} % Thick table middle lines

% GRAMMAR
\usepackage[rounded]{syntax}
\setlength{\grammarparsep}{4pt plus 1pt minus 1pt}

\usepackage{amsmath}
\usepackage{amssymb}

\usepackage{multicol}
\setlength{\columnsep}{1cm}

\usepackage{subcaption}
\usepackage{mathtools}

%\usepackage{floats}
\usepackage[section,above]{placeins}

% Graphs
\usepackage{tkz-berge}
\usetikzlibrary {positioning}
\definecolor {processblue}{cmyk}{0.96,0,0,0}

% Source code
\definecolor{bluekeywords}{rgb}{0.13,0.13,1}
\definecolor{greencomments}{rgb}{0,0.5,0}
\definecolor{redstrings}{rgb}{0.9,0,0}

\usepackage{listings}
\lstdefinestyle{sharpc}{language=[Sharp]C,
	showspaces=false,
	showtabs=false,
	breaklines=true,
	showstringspaces=false,
	breakatwhitespace=true,
	escapeinside={(*@}{@*)},
	commentstyle=\color{greencomments},
	keywordstyle=\color{bluekeywords}\bfseries,
	stringstyle=\color{redstrings},
	basicstyle=\ttfamily,
	morekeywords={var}
	}
%%% STYLE OF PAGES NUMBERING
%\pagestyle{companion} %\nouppercaseheads 
%\pagestyle{headings}
%\pagestyle{ruled}
%\pagestyle{plain}
%\makepagestyle{plain}
% based on companion style
\makepagestyle{hoepelman}
%\newlength{\headwidth}
\setlength{\headwidth}{\textwidth}
  %\addtolength{\headwidth}{\marginparsep}
  %\addtolength{\headwidth}{\marginparwidth}
\makerunningwidth{hoepelman}{\headwidth}
\makeheadrule{hoepelman}{\headwidth}{\normalrulethickness}
\makeheadposition{hoepelman}{flushright}{flushleft}{}{}
\makepsmarks{hoepelman}{%
  \def\chaptermark##1{\markboth{##1}{##1}}    % left mark & right marks
  \def\sectionmark##1{\markright{%
    \ifnum \c@secnumdepth>\z@
      \thesection. \ %
    \fi
    ##1}}
  \def\tocmark{\markboth{\contentsname}{\contentsname}}
  \def\lofmark{\markboth{\listfigurename}{\listfigurename}}
  \def\lotmark{\markboth{\listtablename}{\listtablename}}
  \def\bibmark{\markboth{\bibname}{\bibname}}
  \def\indexmark{\markboth{\indexname}{\indexname}}}
\makepsmarks{hoepelman}{%
  \nouppercaseheads
  \createmark{chapter}{both}{nonumber}{}{}
  \createmark{section}{right}{shownumber}{}{. \space}
  \createplainmark{toc}{both}{\contentsname}
  \createplainmark{lof}{both}{\listfigurename}
  \createplainmark{lot}{both}{\listtablename}
  \createplainmark{bib}{both}{\bibname}
  \createplainmark{index}{both}{\indexname}
  \createplainmark{glossary}{both}{\glossaryname}}
%\makeevenfoot{hoepelman}{}{}{}
%\makeoddfoot{hoepelman}{}{}{}
\makeevenhead{hoepelman}{\normalfont\bfseries\thepage}{}{\normalfont\bfseries\leftmark}
\makeoddhead{hoepelman}{\normalfont\bfseries\thepage}{}{\normalfont\bfseries\leftmark}
%\makeoddhead{hoepelman}{\normalfont\bfseries\rightmark}{}{\normalfont\bfseries\thepage}

\pagestyle{hoepelman}

%%% CHAPTER'S STYLE
\makeatletter
\newcommand\thickhrulefill{\leavevmode \leaders \hrule height 1ex \hfill \kern \z@}
\setlength\midchapskip{10pt}
\makechapterstyle{VZ14}{
\renewcommand\chapternamenum{}
\renewcommand\printchaptername{}
\renewcommand\chapnamefont{\Large\scshape}
\renewcommand\printchapternum{%
\chapnamefont\null\thickhrulefill\quad
\@chapapp\space\thechapter\quad\thickhrulefill}
\renewcommand\printchapternonum{%
\par\thickhrulefill\par\vskip\midchapskip
\hrule\vskip\midchapskip
}
\renewcommand\chaptitlefont{\Huge\scshape\centering}
\renewcommand\afterchapternum{%
\par\nobreak\vskip\midchapskip\hrule\vskip\midchapskip}
\renewcommand\afterchaptertitle{%
\par\vskip\midchapskip\hrule\nobreak\vskip\afterchapskip}
}
\makeatother
\chapterstyle{VZ14}


%%% STYLE OF SECTIONS, SUBSECTIONS, AND SUBSUBSECTIONS
\setsecheadstyle{\Large\bfseries\sffamily\raggedright}
\setsubsecheadstyle{\large\bfseries\sffamily\raggedright}
\setsubsubsecheadstyle{\bfseries\sffamily\raggedright}


% Custom commands
\newcommand\textlcsc[1]{\textsc{\MakeLowercase{#1}}}

\newcommand{\tablelist}[1]{%
%\tightlist%
\begin{itemize}[nosep]%
#1%
\end{itemize}%
\vspace{-\baselineskip}\mbox{}}
%\usepackage[utf8]{inputenc}
\usepackage[T1]{fontenc}
\usepackage{microtype}
%\usepackage{fontspec}

%Default: computer modern

% 
%\usepackage{kpfonts}

% Times-like
%\usepackage{newtxtext}
%\usepackage{newtxmath}

% Times-like
%\usepackage{mathptmx}
%\usepackage{courier}

%Calibri
%\setmainfont{Calibri}

% Palatino
%\usepackage{newpxtext}
%\usepackage{newpxmath}

%\usepackage{libertine}
%\usepackage[libertine]{newtxmath}

\usepackage[scaled]{helvet}
\renewcommand\familydefault{\sfdefault}
\usepackage{sansmath}

\maxsecnumdepth{subsection} % chapters, sections, and subsections are numbered
\maxtocdepth{subsection} % chapters, sections, and subsections are in the Table of Contents

\newcommand{\todo}[1]{\textbf{TODO: #1}}
\newcommand{\ignore}[1]{}
\newcommand{\f}[1]{\texttt{#1}}
\newcommand{\key}[1]{\textbf{#1}}

\begin{document}

\thispagestyle{empty}
\onecolumn
{%%%
\sffamily
\centering

~\vspace{\fill}

{\huge \bfseries
Inter-cell Spreadsheet Refactoring and Parsing Spreadsheet Formulas
}

\vspace{2.0cm}

by

\vspace{2.0cm}

{\LARGE
David Jonathan Hoepelman
}

\vspace{3.0cm}

\textlcsc{Master Thesis} \\
\today

\vspace{2.5cm}

For obtaining the degree of \\
Master of Science \\
in Computer Science - Software Technology \\

\vspace{0.5cm}

Faculty Electrical Engineering, Mathematics and Computer Science (EEMCS)\\
Delft University of Technology

\vspace{1.5cm}

\includegraphics{tudelft}
\hspace{0.5cm}
%\includegraphics[height=8mm]{serg}
%\hspace{0.5cm}
\includegraphics[height=8mm]{spreadsheet-lab}
\hspace{0.5cm}

\vspace{\fill}

}

\cleardoublepage

%%%---%%%---%%%---%%%---%%%---%%%---%%%---%%%---%%%---%%%---%%%---%%%---%%%
\tableofcontents*

\clearpage
%\twocolumn

%%%---%%%---%%%---%%%---%%%---%%%---%%%---%%%---%%%---%%%---%%%---%%%---%%%
%%%---%%%---%%%---%%%---%%%---%%%---%%%---%%%---%%%---%%%---%%%---%%%---%%%

\chapter{Introduction}
Blaatverhaal over waarom spreadsheets veel gebruikt worden en daarom belangrijk zijn, maintainability nightmare en hoe refactoring kan helpen.

\chapter{Spreadsheet anatomy}

\todo{Notificatie dat dit hoofdstuk deels al gepubliceerd is (niet hopen)}

Spreadsheets are cell-oriented dataflow programs.
The spreadsheet model is Turing complete, as proven by an Excel 2010 implementation of a Turing machine \cite{ExcelTuringComplete}.

Altough exotic forms of spreadsheets are available and have been researched, all currently widely used implementation use the following model:
\begin{itemize}
\item A single spreadsheet \key{file} corresponds to a single (\key{work})\key{book}.
\item A workbook can contain any number of (\key{work})\key{sheets}.
\item A sheet consists of a \key{two-dimensional grid} (table) of \key{cells}.
\item A vertical unit in the grid is called a \key{column} and a horizontal unit a \key{row}.
Rows are numbered sequentially top-to-bottom starting at 1, while columns are numbered left-to-right alphabetically, i.e. base-26 using A to Z as digits.
A column or row can also mean all cells contained in that column or row.
\item A \key{cell} can contain a \key{constant value} of any type, a calculation called a \key{formula} or a matrix calculation called an \key{array formula}.
\end{itemize}

\section{Formulas}

Formulas consist of expressions which can contain constant values, \key{function calls} and \key{operators} and, most importantly, \key{references} to other cells.
A cell is identified as a formula cell because all formulas must start with the equals sign \f{=}.

\subsection{Function calls and operators}

Function calls are done by starting with the function name, followed by the arguments in parentheses, separated by a comma.
All spreadsheet implementations provide a range of built-in functions, and in most spreadsheet implementations it is possible to define new functions yourself.
However this is not done directly inside the spreadsheet, and uses an alternate programming language.
In Microsoft Excel and Openoffice.org/Libreoffice this is done with a variant of the BASIC programming language, while in Google Docs this is done with Javascript.

The operators \f{+ - * / = <> <= < >} and \f{>=} can be used according to their usual semantics.
\f{+} and \f{-} are available both unary (\f{=-1}) and binary (\f{=1-1}).
Additionally the \f{\%} unary operator is defined to transform a number into a percentage, \f{\^} is the exponentiation operator and \f{\&} is the text concatenation operator.

Spreadsheets contain three fairly unique binary operators with the range operator \f{:}, the union operator \f{,} and the intersection operator \f{\char32}.
Te semantics of these operators are detailed in section \ref{sec:references}.

\subsection{References}
\label{sec:references}
References are the core component of spreadsheets.
The value of any cell can be used in a formula by concatenating its column and row number, producing a reference like \f{B5}.
This is called A1-style referencing and is by far the most comment in modern spreadsheet implementations.
If the value of a cell changes this new value will be propagated to all formulas that use it.

When copying a cell to another cell by default references will be adjusted by the offset, for example copying \f{=A1} from cell B1 to C2 will cause the copied formula to become \f{=B2}.
This can be prevented by making the reference absolute by prepending a \f{\$} to the column index, row index or both.
The formula \f{=\$A\$1} will remain the same on copy while \f{=\$A1} will still have its row number adjusted when copied.

References can also be \emph{ranges}, which are collections of cells.
Ranges can be constructed by three operators: the range operator \f{:}, the union operator \f{,} (a comma) and the intersection operator \f{\char32} (an amount of whitespace).
The range operator \f{:} creates a rectangular range with the two cells as top-left and bottom-right corners, so \f{=SUM(A1:B10)} will sum all cells in columns A and B with row number 1 through 10.
The range operator is also used to construct ranges of whole rows or columns, for example \f{3:5} is the range of the complete rows three through five, and \f{A:D} is the range of columns A through D.
The union operator, which is different from the mathematical union as duplicates are allowed, combines two references, so \f{A1,C5} will be a range of two cells, \f{A1} and \f{C5}.
Lastly the intersection operator takes only the cells which are in both arguments, \f{=A:A 5:5} will thus be equivalent to \f{=A5}.

A user can also give a name to any collection of cells, thus creating a \emph{named range} which can be referenced in formulas by name.

\subsection{R1C1 reference style}

An alternate style called R1C1 as opposed to the above A1 style exists, but it is very rarely seen or used by users.
In R1C1 reference style one specifies either the offset to a cell between square brackets or its concrete location.
In R1C1 style \f{R[4]C[-2]} means the cell two columns to the left and four rows down, while \f{R2C2} refers to cell B2.
The biggest advantage of R1C1 is that it causes identical formulas to be the same even when they operate on different cells or data because of their position.
These properties make R1C1 useful as an internal representation in spreadsheet implementations.

\subsection{Inter-sheet and external references}
\label{subsection:ExternalRefsDDE}

By default references are to cells or ranges in the same sheet as the formula, but this can be modified with a prefix. A prefix consists of some identifier, followed by an exclamation mark followed by the actual reference.

The most common use case is to reference another sheet in the same workbook, where the prefix is simply the sheet name: \f{=Sheetname!A1}.
References to external spreadsheet files are also possible, which is done by prepending the file name in between square brackets: \f{=[Filename]Sheetname!A1} or \f{=[Filename]!NamedRange}.
A peculiar type of prefix are those that indicate multiple sheets: \f{=Sheet1:Sheet10!A1} means A1 in Sheet1 through Sheet10.
Sheet names can also be between single quotes: \f{='Sheetname with space'!A1}. 

In Windows versions of Microsoft Excel formulas can also call external programs through Dynamic Data Exchange (DDE).
DDE links are a special case of references, used for receiving data from other applications.
They take the form of \f{=Program|Topic!Arguments}, e.g. \f{=Database|TableA!Column1}.

\subsection{Case sensitivity}

Formulas are case-insensitive outside string literals.
Identifiers have a canonical capitalization, and while user can type the identifier with any casing only the canonical form will be displayed.
While the canonical capitalization of built-in identifiers,functions and reserved names, is usually uppercase, the canonical capitalization of an user-defined identifier, an user defined function or a named range, is as the user defined it originally.

\subsection{Whitespace sensitivity}

Outside of string literals formulas are whitespace independent except as a token separator, e.g. \f{=A B} is a formula that uses named ranges \f{A} and \f{B} while the formula \f{=AB} uses only the named range \f{AB}.

An important quirk in this the intersection operator, which acts like a binary operator with the operator being whitespace.
Thus the formula \f{=A\char32B} is the intersection of \f{A} and \f{B}, and so is \f{=A\char32\char32\char32\char32B}.

\section{Array Formulas and Arrays}
\label{sec:arrayformulas}
In spreadsheet programs it is possible to work with one- or two-dimensional matrices.

When constructed from constant values they are called \emph{array constants}, e.g. \f{\{1,2;3,4\}}.
They are surrounded by curly brackets, columns are separated by commas, and rows by semicolons.
Several matrix operations are available, for example \f{=SUM(\{1,2,3\}*10)} will evaluate to 60.

\emph{Array Formulas} use the same syntax as normal formulas, except that the user must enter \emph{Ctrl} + \emph{Shift} + \emph{Enter} to signal that it is an Array formula.
Excel and LibreOffice surround the formula with curly braces.
Google docs works differently and requires the user to surround an array formula with \f{ARRAYFORMULA($\ldots$)}.

Marking a formula as an array formula will enable one- or two-dimensional reference ranges to be treated as matrices, and several matrix operators and functions will be available and allow a formula to return multiple values, to be displayed in multiple cells.
For example if \f{A1},\f{A3},\f{A3} contain the values \f{1},\f{2},\f{3} the array formula \f{\{=SUM(A1:A3*10)\}} will evaluate to \f{60}.
Furthermore, an array formula allows the user to return multiple results, which will be presented in multiple cells.
The array formula \f{\{=\{1,2,3\}*\{4,5,6\}\}} will show \f{4}, \f{10} and \f{18} in three different cells.

While array formulas are powerful, they are also considered complex and not widely used.  \cite{Something}

\section{Type system}

The type system in Microsoft Excel is weak with most typed being able to be coerced into others.
Explicit type conversion is done by certain built-in functions.
For values inside formulas, the following types exist:

\begin{itemize}
\item[Boolean] values are either \f{TRUE} or \f{FALSE}. Booleans can be coerced to numbers where \f{TRUE} will become 1 and \f{FALSE} will become 0 and strings.
\item[Numeric] values are in the range of 8-byte IEEE doubles. Numbers can be provided as integers, decimals or in scientific notation.
When coerced to booleans 0 will become \f{FALSE}, all other values will be \f{TRUE}. Number can also be coerced into strings.
\item[String] values are any Unicode character enclosed in quotation marks \f{"}.
Two quotation marks serve as the escape character, thus \f{""""} represent the string ".
If the contents of a cell start with a \texttt{'} the rest of that cell content is interpreted as a string.

When coerced to booleans all strings except the empty string are \f{TRUE}, the empty string is \f{FALSE}.
When coerced to a numeric value the spreadsheet program will accept any string representing valid numeric user input and otherwise give the error \f{\#VALUE!}.
\item[Error] values are \f{\#DIV/0!}, \f{\#NAME?}, \f{\#NULL!}, \f{\#NUM!}, \f{\#N/A!}, \f{\#VALUE!} and \f{\#REF!}.
Errors behave similar to exceptions or the maybe monad in that they will propagate throughout a calculation.

Errors cannot be coerced, but can be explicitly converted by functions such as \f{ISERROR}.

\item[Ranges and arrays] are one- or two-dimensional matrices of any non-array values. Arrays are rarely used outside of array formulas, but ranges are very common in formulas.
However, these types usually only serve as inputs for functions and are thus fairly transparent to the user outside of array formulas.
Both types usually cannot generally be, doing so will result in the \f{\#VALUE!} error.
\end{itemize}

Some other "display types" exists, these can change the way the data is presented to or validated from the user and can have implications when inter-operating with other programs.
Usually the user can mark a cell as containing one of these types, or Excel can automatically mark a cell to be of this type based on heuristics.
In formulas and internally these types are all represented by one of the above types.
A few of these are commonly used:
\begin{itemize}
\item[Dates and times] are internally stored as a floating point with the integer portion being the number of days since an the epoch January 1st 1900, incorrectly considering 1900 a leap year, and the remainder being the portion of the day that passed.
Excel displays dates and times as is customary in the locale of the user.
When interoperating a date or time value will be exported as a datetime type value of that system.
\item[Currency] is stored as other numbers and displayed in the format customary for the specific currency. When interoperating with some other systems currency values will be exported using arbitrary-precision arithmetic formats.
\item[Percentages] are stored as other numbers, but displayed as if multiplied by 100\%.
\end{itemize}

\chapter{Previous and related work}

\section{Refactoring}

\section{The BumbleBee spreadsheet refactoring suite}

\section{Other spreadsheet refactoring efforts}


\chapter{Candidate spreadsheet refactorings}

\section{Comparison to other programming paradigms}

In the dataflow programming model a calculation is a directed graph where data flows between calculation nodes along the edges.
In a spreadsheet program, nodes are represented by cells which contain a calculation and references to other cells which represent an edge.

Most 

\section{Translating refactorings to the spreadsheet domain}

Tabel met mogelijke refactorings.
Subsecties voor alle geïmplementeerde refactorings.

\chapter{Implementation}

\section{Parsing formulas}

\todo{Notificatie dat deze sectie deels al gepubliceerd is (niet hopen)}

The formula parser component of Bumblebee is called XLParser which is the reference implementation of the grammar by by Aivaloglou, Hoepelman and Hermans \cite{xlparser}.
This component is fully independent, open-source and freely available \footnote{https://github.com/PerfectXL/XLParser}.

XLParser takes a formula and returns a parse tree.
While technically not strictly an AST as it retains several lexical elements, it is also not .
We will indicate the tree returned by XLParser as an AST from here on.

XLParser is built using the Irony Parser framework \footnote{https://irony.codeplex.com}, which uses a grammar defined in C\# to produce a parser based on the LALR(1) algorithm and .NET regular expressions for tokens.

The parser has evolved from previous research efforts, and is intended to be used for ongoing and future research into Microsoft Excel spreadsheet formulas.
With this intent in mind, the following design goals were formulated:

\begin{enumerate}
\label{sec:designgoals}
\item The parser must be compatible with the official language
\item Produced parse trees must be suited for further manipulation and analysis with minimal post-processing required
\item The grammar must be compact enough to feasibly implement with a parser generator
\end{enumerate}

While an official grammar for Excel formulas is published \cite{ExcelOfficialGrammar}, it does not meet the above requirements for two reasons.
Firstly, it is over 30 pages long and contains hundreds of production rules and thus fails requirement 3.
Secondly, because of the detail of the grammar and the large number of production rules the resulting parse trees are very complex and fail requirement 2.

\subsection{Lexical Analysis}
\label{sec:lexanalysis}

\begin{table}
\tiny
\caption{Lexical tokens used in the XLParser grammar, as refered to in section \ref{sec:lexanalysis}.}
\label{table:tokens}
\centerfloat
%\advance\leftskip-1cm
\begin{tabular}{@{}llll@{}}
\toprule
Token Name & Description & Regular Expression & Priority \\
\midrule
%BINOP & Binary Operator & + $\mid$ - $\mid$ / $\mid$ * $\mid$ \textasciicircum $\mid$ \textless $\mid$ \textgreater $\mid$ = $\mid$ \textless= $\mid$ \textgreater= $\mid$ \textless \textgreater & 0 & + \\
BOOL & Boolean literal & TRUE $\mid$ FALSE & 0 \\
CELL & Cell reference & \$? [A-Z]+ \$? [0-9]+ & 2 \\
DDECALL & Dynamic Data Exchange & ' ([A-Z0-9\_ !@\#\$\%\textasciicircum{}\&*()\-+=\{\}:;$\mid$\textless\textgreater,./?\textbackslash{}\textbackslash{}] $\mid$ '')+ ' & 0 \\
ERROR & Error literal & \#NULL! $\mid$ \#DIV/0! $\mid$ \#VALUE! $\mid$ \#NAME? $\mid$ \#NUM! $\mid$ \#N/A & 0 \\
ERROR-REF & Reference error literal & \#REF! & 0 \\
FILE & External file reference & \textbackslash[ [0-9]+ \textbackslash] & 5 \\
FUNCTION & Excel built-in function & (Any entry from the function list\tablefootnote{A function list is available as part of the reference implementation. Lists provide by Microsoft are also available in \cite{ExcelFunctionReference} and \cite{ExcelOfficialGrammar}.}) \textbackslash( & 5        \\
HORIZONTAL-RANGE & Range of rows & \$? [0-9]+ : \$? [0-9]+ & 0 \\
MULTIPLE-SHEETS & Multiple sheet references & \begin{tabular}[c]{@{}l@{}}
[A-Z0-9]+ : ([A-Z0-9\_.]+$\mid$'([A-Z0-9\_ !\textasciicircum{}\&*()\-+=\{\}:;$\mid$\textless\textgreater,./?\textbackslash{}\textbackslash{}] $\mid$ '')+')!
\end{tabular}
& 1 \\
NAMED-RANGE & Named range & [A-Z\_][A-Z0-9\_.]* & -2 \\
NAMED-RANGE-PREFIXED & \begin{tabular}[c]{@{}l@{}} Named range which starts with \\ a string  that could be another token \end{tabular} & (TRUE $\mid$ FALSE $\mid$ [A-Z]+[0-9]+)    {[}A-Z0-9\_.{]}+                                                                             & 3 \\
NUMBER-LITERAL & \begin{tabular}[c]{@{}l@{}}An integer, floating point\\     or scientific notation number literal\end{tabular} & [0-9]+ ,? [0-9]* (e [0-9]+)? & 0 \\
QUOTED-FILE-SHEET & A file reference within single quotes & '\textbackslash[ [0-9]+ \textbackslash] ([0-9A-Z\_ !@\#\$\%\textasciicircum{}\&*()\-+=\{\}:;$\mid$\textless\textgreater,./?\textbackslash{}\textbackslash{}] $\mid$ '')+ '!
& 5        \\
REFERENCE-FUNCTION & Excel built-in reference function & (INDEX $\mid$ OFFSET $\mid$ INDIRECT)\textbackslash( & 5 \\
RESERVED-NAME & An Excel reserved name & \_xlnm\textbackslash.  [A-Z\_]+ & -1 \\
SHEET & The name of a worksheet &
	([0-9A-Z\_.]+ $\mid$ ' ([0-9A-Z\_ !@\#\$\%\textasciicircum{}\&*()\-+={}$\mid$:;\textless\textgreater,./?\textbackslash\textbackslash] $\mid$ '')+ ') !
& 5        \\
STRING & String literal & " ([\textasciicircum{} "] $\mid$ "")* " & 0       \\
%UNOP\_POSTFIX & Unary postfix operator & \% & 0 & \% \\
%UNOP\_PREFIX & Unary prefix operator & + $\mid$ - & 0 & -                  \\
UDF & User Defined Function & (\_xll\textbackslash.)? [A-Z0-9]+  ( & 4 \\
VERTICAL-RANGE & Range of columns & \$? [A-Z]+ : \$? [A-Z]+ & 0 \\ 
\bottomrule
\end{tabular}
\end{table}

\todo{Move table to appendix}

Table \ref{table:tokens} contains the lexical tokens of the grammar, along with their identification patterns in a simple regular expression language. All tokens are case-insensitive.

This token list requires the scanner to support token priorities, which the Irony framework does.
While it is possible to remove this necessity by altering the tokens and production rules, this would make the resulting grammar and AST harder to use.

Some simple tokens are directly defined in the production rules in Figure \ref{figure:productions} in between quotes for readability and compactness.

\subsection{Syntactical Analysis}

\begin{figure}
\small
\caption{Production rules}
\label{figure:productions}
\begin{multicols*}{2}
\begin{grammar}
<Start> ::= <Constant>
	\alt '=' <Formula>
	\alt <ArrayFormula>
	
<ArrayFormula> ::= `\{=' <Formula> `\}'

<Formula> ::= <Constant>
         \alt <Reference>
         \alt <FunctionCall>
         \alt `(' <Formula> `)'
         \alt <ConstantArray>
         \alt "RESERVED-NAME"
         
<Constant> ::= "NUMBER" | "STRING" | "BOOL" | "ERROR"
%         \alt "STRING"
%         \alt "BOOL"
%         \alt "ERROR"
         
<FunctionCall> ::= <Function> <Arguments> `)'
		\alt <UnOpPrefix> <Formula>
		\alt <Formula> `\%'
		\alt <Formula> <BinOp> <Formula>
		
<UnOpPrefix> = `+' | `-'

<BinOp> = `+' | `-' | `*' | `/' | `\textasciicircum'
	\alt `<' | `>' | `=' | `<=' | `>=' | `<>'

<Function> ::= "FUNCTION" | "UDF"
	
<Arguments> ::= <Argument>
	\alt <Argument> `,' <Arguments>

<Argument> ::= <Formula> | $\epsilon$

<Prefix> ::= "SHEET"
	\alt "FILE" "SHEET"
	\alt "FILE" `!'
	\alt "QUOTED-FILE-SHEET"
	\alt "MULTIPLE-SHEETS"
	\alt "FILE" "MULTIPLE-SHEETS"

\columnbreak

<Reference> ::= <ReferenceItem>
	\alt <Reference> `:' <Reference>
	\alt <Reference> `\ ' <Reference>
	\alt `(' <Union> `)'
	\alt `(' <Reference> `)' 
	\alt <Prefix> <ReferenceItem>
    \alt <Prefix> "UDF" <Arguments> `)'
    \alt <DynamicDataExchange>
    
<ReferenceItem> ::= "CELL"
	\alt <NamedRange>
	\alt "REFERENCE-FUNCTION" Arguments `)'
	\alt "VERTICAL-RANGE"
	\alt "HORIZONTAL-RANGE"
	\alt "ERROR-REF"
	
<NamedRange> ::= "NAMED-RANGE"
            \alt "NAMED-RANGE-PREFIXED"

<Union> ::= <Reference> | <Reference> `,' <Union>
	
<DynamicDataExchange> ::= "FILE" `!' "DDECALL"

<ConstantArray> ::= `\{' <ArrayColumns> `\}'

<ArrayColumns> ::= <ArrayRows>
	\alt <ArrayRows> `;' <ArrayColumns>

<ArrayRows> ::= <ArrayConstant>
	\alt <ArrayConstant> `,' <ArrayRows>
	
<ArrayConstant> ::= <Constant>
	\alt <UnOpPrefix> "NUMBER"
	\alt "ERROR-REF"

\end{grammar}
\end{multicols*}
\end{figure}

The complete production rules of the used grammar can be found in Figure \ref{figure:productions} in Extended BNF syntax.
The start symbol is $Start$.

\synt{Formula} and \synt{Reference} are the two most important production rules in this grammar.
They model respectively all expressions inside formulas, and expressions consisting purely of references.
This distinction is important because in some places in the language only reference expression can be used.
An example of this is the the range operator \f{:} which only accepts reference expressions as operands.
Thus a formula like \f{=1.5:2.5} will result in a parse error in both Microsoft Excel and XLParser, while a formula like \f{="a" + 2} will result in a runtime type system error.

The \synt{Formula} and \synt{Reference} production rules are also illustrated as syntax diagrams, with most production rules expanded, in Figures \ref{figure:Formula} and \ref{figure:Reference}.

\todo{Move these figures to appendix?}

\begin{figure}
	\caption{Syntax diagram of the \synt{Formula} production rule with most production rules expanded}
	\label{figure:Formula}
	%\centerfloat
	\begin{grammar}
	<Formula> ::= \begin{syntdiag}
	\begin{stack} 
	\begin{stack}
	\begin{stack} "FUNCTION" \\ "UDF" \end{stack} \begin{rep} \begin{stack} \\ <Formula> \end{stack} \\  '$,$' \end{rep} '$)$' \\
	\begin{stack} \\ '$+$' \\ '$-$' \end{stack} <Formula> \begin{stack} \\ '$\%$' \end{stack}\\ 
	%<Formula>  \begin{stack} '$+$' \\ '$-$' \\ '$*$' \\ '$/$' \\ '\textasciicircum'  \\ '\textless' \\ '\textgreater' \\ '$=$' \\ '$\textless=$' \\ '$\textgreater=$' \\ '$\textless\textgreater$'   \end{stack} <Formula>
	<Formula> <BinOp> <Formula>
	\end{stack} \\
	\begin{stack} "NUMBER" \\ "STRING" \\ "BOOL" \\ "ERROR"  \end{stack} \\ "RESERVED-NAME" \\ <Reference> \\ '$($' <Formula> '$)$' \\  '$\{$' \begin{rep} \begin{rep}	\begin{stack} \begin{stack} \begin{stack} \\ '$+$' \\ '$-$' \end{stack} "NUMBER" \\ "STRING" \\ "BOOL" \\ "ERROR"  \end{stack} \\ "ERROR-REF" \end{stack} \\  '$,$' \end{rep} \\ '$;$' \end{rep} '$\}$'
	\end{stack}
	\end{syntdiag}
\end{grammar}
\end{figure}

\begin{figure}
	\caption{Syntax diagram of the \synt{Reference} production rule with most production rules expanded}
	\label{figure:Reference}
	\begin{grammar}
		<Reference> ::= \begin{syntdiag}[\scriptsize\sdlengths]
		\begin{stack} '$($' \begin{rep} <Reference> \\  '$,$' \end{rep} '$)$'\\ <Reference> \begin{stack} '$:$' \\ '\char32' \end{stack} <Reference> \\
		\begin{stack} \\ \begin{stack} \\ "FILE" \end{stack} \begin{stack} "SHEET" \\ "MULTIPLE_SHEETS" \end{stack} \\ "FILE" '$!$' \\ "QUOTED_FILE_SHEET" \end{stack}
		\begin{stack} \begin{stack} "CELL" \\ "VERTICAL-RANGE" \\ "HORIZONTAL-RANGE" \\ \begin{stack} "NAMED-RANGE" \\ "NAMED-RANGE-COMBINED" \end{stack} \\ "ERROR-REF" \end{stack} \\  \begin{stack} "REFERENCE-FUNCTION" \\ "UDF" \end{stack} \begin{rep} \begin{stack} \\ <Formula> \end{stack} \\  '$,$' \end{rep} '$)$'\end{stack}
		\\"FILE" '$!$' "DDECALL"
		\end{stack}
		\end{syntdiag}
	\end{grammar}
\end{figure}

\subsection{Precedence and ambiguity}

\begin{table}
\small
\caption{Operator precedence in formulas}
\label{table:operatorprec}
\begin{tabular}{lll}
Precedence & Operator(s) & Associativity \\
1 & = \textless \  \textgreater \  \textless= \  \textgreater= \  \textless\textgreater & Left         \\
2 & \& & Left \\
3 & + - (binary) & Left \\
4 & $\ast$ / & Left \\
5 & \textasciicircum & Left\footnotemark \\
6 & \% & \\
7 & + - (unary) & \\
8 & : , \texttt{\char32} & Left
\end{tabular}
\end{table}

\footnotetext{This is contrary to most other languages, where the exponentiation operator is right-associative. \\
In Excel \f{2\textasciicircum 1\textasciicircum 2} will be $(2^1)^2 = 4$, while in most other languages it will be $2^{1^2} = 2$}

The production rules are ambiguous, which means they cannot be directly used in a parser generator based on the LALR algorithm like Irony.

To resolve ambiguity with operators, e.g. whether to parse \f{=1+2*3} as \f{=(1+2)*3} or \f{=1+(2*3)}, operator precedence and associativity rules are defined.
These can be found in table \ref{table:operatorprec}.

However, even with precedence and associativity rules the grammar is still not fully un-ambigious.
This is due to trade-offs on parsing references, see section \ref{tradeoff:references}, and parsing unions (see section \ref{subsec:desing:unions}).
Ambiguity exists between the following production rules:
\begin{enumerate}
\item \begin{grammar}<Reference> ::= `(' <Reference> `)'\end{grammar}
\item \begin{grammar}<Reference> ::= `(' <Union> `)'\end{grammar}
\item \begin{grammar}<Formula> ::= `(' <Formula> `)'\end{grammar}
\end{enumerate}

A formula like \texttt{=(A1)} can be interpreted as either a bracketed reference, a union of one reference, or a reference within a bracketed formula.

In a LALR like the one Irony produces this ambiguity manifests in a state where, on a \texttt{')'} token, shifting on rule 1 and reducing on either rule 2 or 3 are possibilities, causing a shift-reduce conflict.
This was solved by instructing the parser generator to shift on rule 1 in case of this conflict, because this always is a correct interpretation and thus results in correct ASTs.

\subsection{Trade-offs and quirks}

The grammar presented in this chapter (\todo{section?}) contains some quirks on sub-optimalities, partly due to the Excel language itself, partly due to design decisions.
These are detailed here

\subsubsection{\textbf{References}}
\label{tradeoff:references}

In some places only reference expressions are accepted, instead of all expressions.
Furthermore references are of great importance in spreadsheet formulas, and thus of interest for analysis.
To support easier analysis (design goal 2) references have different production rules than other expressions.
This causes references to be easily identified and isolated, but has the downside of increasing ambiguity, as explained in Section \ref{sec:ambiguity}. 

Another approach would be to parse all expressions with the same production rules and implement a type system, however this would be very detrimental to both ease of analysis (design goal 2) and to ease of implementation (design goal 3).

\subsubsection{\textbf{Unions}}
\label{subsec:desing:unions}

The comma serves both as an union operator and a function argument separator.
This proves challenging to correctly implement in a LALR(1) grammar.

A straightforward implementation would use production rules similar to this:
\begin{grammar}
<Union> ::= <Reference> `,' <Reference>

<Arguments> ::= <Argument>
	\alt <Argument> `,' <Arguments>
\end{grammar}

However, this will cause a reduce-reduce conflict because the parser will have a state wherein it can reduce to both a \synt{Union} or \synt{Argument} on a \texttt{,} token.
Unfortunately there is no correct choice: in a formula like \texttt{=SUM(A1,1)} the parser must reduce on the \synt{Argument} nonterminal, while in a formula like \texttt{=A1,A1} the parser must reduce to the \synt{Union} nonterminal.
With the above production rules a LALR(1) parser could not correctly parse the language.

The presented grammar only parsers unions in between parentheses, e.g. \texttt{=SMALL((A1,A2),1)}.
This is a trade-off between a lower compatibility (design goal 1) and an easier implementation (design goal 3).
Because unions are very rare \cite{grammarpaper} this lower compatibility is deemed acceptable.

Additionally formulas that this grammar does not parse often result in runtime errors after evaluation.
For example \texttt{=A1,A1} does parse in a spreadsheet program, but produces the error \texttt{\#VALUE!} on evaluation.

Implementing the straightforward rules above, while desirable, is not possible without using a more powerful grammar class.


\subsection{Trade-offs and quirks in the grammar}

\subsection{Excel Formula Grammar}

\section{Transforming formula ASTs}

\section{Implementing refactorings}

\subsection{Extract formula}

\subsection{Inline formula}

\subsection{Addition to SUMIF}

\subsection{\textasciicircum\textasciicircum or Aggregate to conditional Aggregate}

\subsection{Group References}

\chapter{Evaluation}


\chapter{Conclusion}

\section{Future Work}

\bibliographystyle{unsrt}
\bibliography{thesis}

\end{document}

